\section{Total degree: Schwartz--Zippel, lines, and the application}

\begin{lemma}[Schwartz--Zippel for the total degree]\label{lem:sz-total}
Let $m\ge1$, $D\in\N$, and let $f\in\F[x_1,\dots,x_m]$ be nonzero with
$\deg_{\mathbf 1}f\le D$, where $\deg_{\mathbf 1}:=\deg_w$ for $w=\mathbf 1$ (\cref{def:wdeg}). For a
uniform $r\in\F^m$,
\[
\Pr\big[f(r)=0\big]\le\frac{D}{|\F|}.
\]
From the individual bounds alone, that is from $\deg_{x_l}f<k$ for all $l$ (the box
$\F[x_1,\dots,x_m]_{<k}$ of \cref{def:poly}), one gets $\deg_{\mathbf 1}f\le m(k-1)$ and hence the
bound $m(k-1)/|\F|$; this is the bound used in \cref{prop:complete}(c).
\end{lemma}

\begin{proof}\begin{enumerate}
\item Induction on $m$. For $m=1$, a nonzero $f\in\F[T]$ has at most $\deg f\le D$ roots, so the claim holds.
\item For $m\ge2$ write $f=\sum_{j\le d}f_j(x_1,\dots,x_{m-1})x_m^j$ with $f_d\ne0$; if $d=0$ then $f=f_0$ does not involve $x_m$ and the claim is the induction hypothesis, so let $d\ge1$.
\item Then $\deg_{\mathbf 1}f_d+d\le\deg_{\mathbf 1}f$. Write $r=(r',r_m)$ with $r'$ uniform on $\F^{m-1}$; by the induction hypothesis applied to the $(m-1)$-variate polynomial $f_d$, $\Pr[f_d(r')=0]\le\deg_{\mathbf 1}f_d/|\F|\le(\deg_{\mathbf 1}f-d)/|\F|$.
\item Conditioned on $f_d(r')\ne0$ the polynomial $T\mapsto f(r',T)$ is nonzero of degree $d$, so it vanishes at the uniform $r_m$ with probability at most $d/|\F|$.
\item Adding the two bounds of steps 3 and 4 gives $\Pr[f(r)=0]\le(\deg_{\mathbf 1}f-d+d)/|\F|=\deg_{\mathbf 1}f/|\F|\le D/|\F|$.
\end{enumerate}\end{proof}
\uses{\cref{def:wdeg}}\usedby{\cref{rem:total-app}}

\begin{lemma}[Restriction to a line]\label{lem:line-total}
Let $m\ge1$, $f\in\F[x_1,\dots,x_m]$ and $\mathbf a,\mathbf b\in\F^m$. Then
\[
f(\mathbf a+T\mathbf b):=f(a_1+Tb_1,\dots,a_m+Tb_m)\in\F[T]
\]
has degree at most $\deg_{\mathbf 1}f$.
\end{lemma}

\begin{proof}\begin{enumerate}
\item The substitution $x_l\mapsto a_l+Tb_l$ sends $x^\alpha$ to $\prod_{l\le m}(a_l+Tb_l)^{\alpha_l}$, whose $T$-degree is at most $\alpha_1+\dots+\alpha_m$, and to $0$ when $a_l=b_l=0$ for some $l$ with $\alpha_l\ne0$.
\item By \cref{def:wdeg} with $w=\mathbf 1$, $\langle\mathbf 1,\alpha\rangle=\alpha_1+\dots+\alpha_m\le\deg_{\mathbf 1}f$ for every $\alpha$ with $f_\alpha\ne0$, so every term of $f(\mathbf a+T\mathbf b)$ has $T$-degree at most $\deg_{\mathbf 1}f$.
\item A sum of polynomials of $T$-degree at most $\deg_{\mathbf 1}f$ has $T$-degree at most $\deg_{\mathbf 1}f$, and $0$ has $T$-degree $-\infty$.
\end{enumerate}\end{proof}
\uses{\cref{def:wdeg}}\usedby{\cref{rem:total-app}}

\begin{remark}[Two evaluation claims reduced to one]\label{rem:total-app}
Let $f$ with $\deg_{x_l}f<k$ for all $l$ be committed, and suppose that $\Pi_{\mathbf 1,D}$
(\cref{def:protocol}) has certified $\deg_{\mathbf 1}f\le D$. A protocol may then have to convince
the verifier of two claims $f(\mathbf a)=y_{\mathbf a}$ and $f(\mathbf b)=y_{\mathbf b}$, for
$\mathbf a,\mathbf b\in\F^m$ and $y_{\mathbf a},y_{\mathbf b}\in\F$. Let
$\ell(T):=\mathbf a+T(\mathbf b-\mathbf a)$, so that $\ell(0)=\mathbf a$ and $\ell(1)=\mathbf b$. The
prover sends a polynomial $q\in\F[T]_{\le D}$ claimed to be $f\circ\ell$; the verifier checks
$q(0)=y_{\mathbf a}$ and $q(1)=y_{\mathbf b}$, draws a uniform $T^*\in\F$ and asks for an opening of
the commitment at $\ell(T^*)\in\F^m$. If $q\ne f\circ\ell$, then $q-f\circ\ell$ is nonzero of degree
at most $D$ by \cref{lem:line-total}, so $q(T^*)\ne f(\ell(T^*))$ with probability at least
$1-D/|\F|$ by \cref{lem:sz-total} with $m=1$. The two claims are thus replaced by one, at error
$D/|\F|$ and at a cost of $D+1$ field elements. With only the individual bounds,
$\deg_{\mathbf 1}f\le m(k-1)$ is all that is known, the same step costs $m(k-1)+1$ field elements and
its error is $m(k-1)/|\F|$. At
$k=4$, $m=8$, $D=6$ the certificate gives $7$ field elements at error $6/|\F|$, against $25$
elements at $24/|\F|$ without it.
\end{remark}
\uses{\cref{lem:sz-total,lem:line-total,def:protocol}}\usedby{the paper}

\begin{remark}[Protocols whose soundness is stated in the total degree]\label{rem:total-prior}
For some interactive proofs the total degree is the parameter of the soundness statement, and a
certificate of $\deg_{\mathbf 1}f\le D$ is then an input to their analysis rather than a
strengthening of it. Fold-DCS [Fold-DCS, arXiv:2504.00693] is one: "the soundness of Fold-DCS
depends on the total degree of the polynomial, not its individual degrees" (p.~2). Its error is
bounded by $\log(m)D/|\F|$ for $m$ variables and total degree $D$, against $mb/|\F|$ for the
classical sumcheck with individual degree bound $b$ (their $\mu$ and $d$) (\S3.2.2), and it constrains the commitment itself
to total degree at most $D$ (\S4.2.1). The line step of \cref{rem:total-app} is the one of the GKR
protocol [GKR, STOC 2008 / J.~ACM 62(4), 2015], which uses it to reduce the two claims of a round to
a single one [Thaler, \S4.6.4]; the subroutine, with its soundness claim, is
[Thaler, \S4.5.2, Claim~4.6].
\end{remark}
\uses{\cref{rem:total-app}}\usedby{the paper}
